\documentclass[a4paper,12pt]{ctexart}
\usepackage{xcolor}
\usepackage[top=1.8cm, bottom=1.8cm, left=1.6cm, right=1.6cm]{geometry}
\usepackage{array}
\usepackage{longtable}
\usepackage{hyperref}
\hypersetup{colorlinks=true,linkcolor=blue!70!black}
\linespread{1.35}

\definecolor{donegreen}{RGB}{40,140,80}
\definecolor{planblue}{RGB}{30,100,180}

\newcommand{\done}{\textcolor{donegreen}{\textbf{[已完成]}}}
\newcommand{\plan}{\textcolor{planblue}{\textbf{[待写]}}}

\title{\textcolor{blue!70!black}{\Huge\textbf{自然科学系列 · 整体规划大纲}}}
\author{}
\date{}

\begin{document}
\maketitle
\thispagestyle{empty}

\section*{\textcolor{blue!70!black}{总体规划思路}}

面向北京地区四年级起点的孩子，覆盖初中（对标北京中考）和高中（对标北京高考3+3选考）的自然科学内容。按"启蒙$\rightarrow$分科入门$\rightarrow$深度进阶$\rightarrow$拓展竞赛"四层递进。

\textbf{学习路径：}
\begin{enumerate}
  \item 启蒙入口：先通读\textbf{自然科学小故事（28课）}，建立直觉和好奇心
  \item 分科基础：进入各科的"入门"本（物理1-3、化学1-2、生物1-2、地理1）
  \item 深度进阶：升入各科的"进阶"本（物理4-6、化学3-4、生物3-4、地理2）
  \item 拓展竞赛：有额外兴趣或竞赛需求时，选读拓展本（可选）
\end{enumerate}

\vspace{0.3cm}
\textcolor{gray}{\small 总计20本（含已完成1本），每本约20-30课。各本之间可交叉阅读。}

\newpage

% ============================================================
\section*{\textcolor{orange!80!black}{一、启蒙模块（1本，已完成）}}
% ============================================================

\vspace{0.2cm}
{\large\textbf{自然科学小故事} \done}
\vspace{0.1cm}

28个故事，不讲公式不背概念。物理/化学/生物/地科各约占四分之一。从苹果落地到板块运动，每个"为什么"引出一个小故事。\\
\textcolor{gray}{文件：science/nature-intro/classroom\_materials.tex}

\vspace{0.5cm}

% ============================================================
\section*{\textcolor{orange!80!black}{二、物理系列（6本）}}
% ============================================================

\subsection*{物理1 · 力学入门 \plan}

\textbf{对标：} 初中八上+八下（中考计分）\\
\textbf{内容：} 机械运动（参照物、速度、$v = s/t$）、力（重力/弹力/摩擦力/二力平衡）、压强（固体/液体/大气）、浮力（$F_{\text{浮}} = \rho g V_{\text{排}}$）、简单机械（杠杆/滑轮）、功与机械效率。\\
\textbf{定位：} 物理的第一本。从"推箱子""游泳感觉轻""跷跷板"等生活经验切入，逐步引入物理量和公式。注重实验测量方法。

\subsection*{物理2 · 声光热 \plan}

\textbf{对标：} 初中八上后半+九上热学\\
\textbf{内容：} 声（产生传播/音调响度音色）、光（直线传播/反射定律/折射定律/凸透镜成像）、热（温度计/物态变化）、内能（比热容$Q = cm\Delta t$/热机效率）。\\
\textbf{定位：} 三种"感官物理"打包一本。凸透镜成像规律是初中物理最大难点之一，需重点着墨。

\subsection*{物理3 · 电磁入门 \plan}

\textbf{对标：} 初中九年级全册\\
\textbf{内容：} 电路（组成/串并联/电流表电压表）、欧姆定律（$I=U/R$）、电功率（$P=UI$/$Q=I^2Rt$）、电磁（电生磁/电磁铁/电动机/发电机）。\\
\textbf{定位：} 初中物理最抽象的一块——电看不见摸不着。需大量电路图和实验示意。中考重头，多设例题。

\subsection*{物理4 · 力学与运动定律 \plan}

\textbf{对标：} 高中必修1\\
\textbf{内容：} 匀变速直线运动（位移/速度/加速度/$v$-$t$图像/自由落体）、牛顿三定律的系统讲解（受力分析/整体法与隔离法/超重与失重）。\\
\textbf{定位：} 高中物理的"地基本"。从初中"定性+简单定量"跃迁到高中"矢量分析+系统受力"。牛顿定律是分水岭——这一本打不牢，后面全摇晃。

\subsection*{物理5 · 曲线、能量与动量 \plan}

\textbf{对标：} 高中必修2+选择性必修1前半\\
\textbf{内容：} 曲线运动（运动的合成与分解/平抛/圆周与向心力）、万有引力与航天（开普勒定律/宇宙速度）、机械能守恒（功与能/动能定理）、动量守恒（动量/冲量/碰撞）。\\
\textbf{定位：} 三大守恒律（能+动量）是高中物理的核心思想。天体运动可呼应启蒙本的牛顿苹果故事。

\subsection*{物理6 · 电磁、波与近代物理 \plan}

\textbf{对标：} 高中必修3+选择性必修1后半+选择性必修2+选择性必修3精选\\
\textbf{内容：} 静电场（场强/电势/电容器）、恒定电流（闭合电路欧姆定律深化）、磁场与电磁感应（安培力/洛伦兹力/法拉第定律/楞次定律）、振动与波（简谐运动/机械波/干涉衍射）、光学（波动光学/几何光学深化）、热力学（分子动理论/热力学定律）、原子与核物理（光电效应/玻尔模型/核反应）、相对论简介（时间膨胀/$E=mc^2$）。\\
\textbf{定位：} 高中物理容量最大的一本。电磁感应是公认最难点。近代物理以科普+基本计算为主。热力学和原子物理可适度精简。

\newpage

% ============================================================
\section*{\textcolor{orange!80!black}{三、化学系列（4本）}}
% ============================================================

\subsection*{化学1 · 物质构成与化学语言 \plan}

\textbf{对标：} 初中九上 + 高中必修1前半\\
\textbf{内容：} 物理变化与化学变化、原子模型演进（道尔顿$\rightarrow$卢瑟福$\rightarrow$玻尔$\rightarrow$电子云）、元素与周期表（周期/族/分区）、化合价与化学式书写、化学方程式的配平（最小公倍数法）、质量守恒定律、物质的量（摩尔/阿伏加德罗常数）、化学键（离子键/共价键/分子空间构型）。\\
\textbf{定位：} 化学的"语言课"——从实物到符号的抽象过程是第一道坎。把化学式当"密码"来学，化学方程式就是"用化学语言写的一句话"。

\subsection*{化学2 · 常见元素与无机反应 \plan}

\textbf{对标：} 初中九上后半+九下 + 高中必修1后半\\
\textbf{内容：} 空气与氧气（组成/性质/制法/燃烧条件）、水与氢气、碳家族（金刚石/石墨/CO/CO$_2$/碳酸钙）、金属（活动性顺序/铁的冶炼/铝的氧化膜）、溶液（溶解度/结晶）、酸碱盐（pH/复分解反应/盐的溶解性）、氧化还原反应（化合价升降法配平）、离子反应（离子方程式书写）、元素化学（钠/氯/硫/氮的化合物体系）。\\
\textbf{定位：} 从具体物质出发建立"物质$\rightarrow$性质$\rightarrow$反应"的认知链。初中偏描述和简单推理，高中引入氧化还原和离子反应的系统分析框架。

\subsection*{化学3 · 化学反应原理 \plan}

\textbf{对标：} 高中选择性必修1\\
\textbf{内容：} 反应热与焓变（热化学方程式/盖斯定律）、化学反应速率（浓度/温度/压强/催化剂）、化学平衡（平衡常数/勒夏特列原理）、水溶液中的离子平衡（弱电解质电离/$K_w$/pH计算/盐类水解/沉淀溶解平衡）、电化学（原电池/电解池/金属腐蚀与防护）。\\
\textbf{定位：} 高中化学最难的一本——四大平衡体系（化学平衡/电离/水解/沉淀）构成核心框架。需大量计算例题和图像分析。

\subsection*{化学4 · 有机化学与物质结构 \plan}

\textbf{对标：} 高中选择性必修2+选择性必修3\\
\textbf{内容：} 有机物的分类/同分异构/命名、烃（烷/烯/炔/芳香烃）、烃的衍生物（卤代烃/醇酚/醛酮/羧酸酯）、生物大分子（糖类/油脂/蛋白质/核酸）、合成高分子（塑料/橡胶/纤维）、原子结构深化（电子层与能级/构造原理）、分子结构（杂化轨道/VSEPR/分子极性）、晶体结构（离子/分子/共价/金属晶体）。\\
\textbf{定位：} 有机化学需要建立"碳骨架$\rightarrow$官能团$\rightarrow$反应类型"的新思维。结构部分抽象，多用3D示意图辅助。

\newpage

% ============================================================
\section*{\textcolor{orange!80!black}{四、生物系列（4本）}}
% ============================================================

\textcolor{gray}{划分逻辑：统一按"生物学组织尺度"递进——从分子级逐层放大到生物圈级。}

\subsection*{生物1 · 分子与细胞 \plan}

\textbf{尺度：} 分子 $\rightarrow$ 细胞\\
\textbf{对标：} 初中七上（细胞部分）+ 高中必修1\\
\textbf{内容：}
\begin{itemize}
  \item 细胞的分子组成（高中）：水、无机盐、糖类、脂质、蛋白质、核酸
  \item 细胞的结构：细胞膜、细胞质、细胞核、主要细胞器（线粒体/叶绿体/内质网/高尔基体等）
  \item 细胞的代谢（高中）：酶的特性、ATP、细胞呼吸（有氧/无氧）、光合作用（光反应/暗反应）
  \item 细胞的增殖：有丝分裂（各时期特征）、减数分裂（高中）
  \item 细胞的分化与衰老：干细胞、细胞凋亡
\end{itemize}
\textbf{定位：} 生物学的"分子级地基"。初中以形态描述为主，高中深入代谢机理。光合作用和呼吸作用是最大难点，能量转换的化学本质需讲透彻。

\subsection*{生物2 · 生物体的结构层次 \plan}

\textbf{尺度：} 细胞 $\rightarrow$ 组织 $\rightarrow$ 器官 $\rightarrow$ 系统 $\rightarrow$ 完整个体\\
\textbf{对标：} 初中七上（结构层次+植物）+ 初中七下（人体）+ 初中八上（动物）\\
\textbf{内容：}
\begin{itemize}
  \item 结构层次总论：细胞分化$\rightarrow$组织$\rightarrow$器官$\rightarrow$系统$\rightarrow$个体
  \item 植物体的结构：根茎叶花果实种子、输导组织（导管/筛管）、植物的生理（光合/呼吸/蒸腾/向性运动）
  \item 动物体的结构（以人为代表）：人体八大系统——消化、呼吸、循环、泌尿、神经、内分泌、运动、生殖——逐一讲解结构和功能
  \item 动物的行为：先天性行为与学习行为
\end{itemize}
\textbf{定位：} "从细胞到完整个体"一条线拉通。植物和动物（人体）两种"构建方案"在同一本里对比。人体系统是七下核心，配大量示意图。

\subsection*{生物3 · 遗传与进化 \plan}

\textbf{尺度：} 个体 $\rightarrow$ 种群（时间维度）\\
\textbf{对标：} 初中八下（生殖+遗传）+ 高中必修2\\
\textbf{内容：}
\begin{itemize}
  \item 生物的生殖：无性生殖与有性生殖
  \item 孟德尔遗传定律：分离定律与自由组合定律、概率计算
  \item 基因与染色体：伴性遗传、基因的连锁与交换（高中）
  \item 基因的本质（高中）：DNA双螺旋结构、半保留复制、转录与翻译（中心法则）
  \item 变异：基因突变、基因重组、染色体变异
  \item 生物的进化：达尔文自然选择学说、现代综合进化论、物种形成
\end{itemize}
\textbf{定位：} 遗传学是高中生物最"理科"的板块——孟德尔定律+概率计算是核心。进化论是贯穿全书的"大故事"，和启蒙本的恐龙灭绝课可呼应。

\subsection*{生物4 · 生态与生物多样性 \plan}

\textbf{尺度：} 种群 $\rightarrow$ 群落 $\rightarrow$ 生态系统 $\rightarrow$ 生物圈\\
\textbf{对标：} 初中八上（微生物+多样性）+ 初中八下（生态）+ 高中选择性必修2\\
\textbf{内容：}
\begin{itemize}
  \item 微生物：细菌/真菌/病毒的结构与繁殖、与人类的关系
  \item 生物分类与多样性：界门纲目科属种、物种多样性/基因多样性/生态系统多样性
  \item 种群生态（高中）：种群数量特征、种群增长模型（J型/S型）
  \item 群落生态：种间关系（捕食/竞争/寄生/互利共生）、群落演替
  \item 生态系统：食物链与食物网、能量流动（十分之一定律）、物质循环（碳循环/氮循环）
  \item 生态保护：生物多样性的价值、自然保护区、人类活动对生态的影响
\end{itemize}
\textbf{定位：} 从微观回到宏观。微生物、分类、多样性这些初中"零散"内容统合在"生态系统"这个大框架下。能量流动和物质循环是高中重点。

\newpage

% ============================================================
\section*{\textcolor{orange!80!black}{五、地理系列（2本）}}
% ============================================================

\textcolor{gray}{说明：本系列仅覆盖自然地理。人文/经济/区域地理划归 social-science/ 系列。}

\subsection*{地理1 · 自然地理基础 \plan}

\textbf{对标：} 初中七上 + 高中必修1\\
\textbf{内容：}
\begin{itemize}
  \item 地球的宇宙环境：地球的形状与大小、自转（昼夜与时区）与公转（四季与五带）、经纬度
  \item 地图：比例尺、方向（指向标/经纬网）、等高线地形图的判读
  \item 大气与气候：气温与降水、世界主要气候类型及分布
  \item 大气运动（高中深化）：热力环流、气压带与风带、季风环流
  \item 水循环（高中深化）：海陆间循环/陆地内循环/海上内循环、洋流（风海流/密度流/补偿流）
  \item 自然环境的整体性与差异性（高中）：纬度地带性、经度地带性、垂直地带性
\end{itemize}
\textbf{定位：} 地理的硬核基础。地球运动和气候成因是初高中地理最大难点——空间思维要求高。等高线判读是终身实用技能。此本可与启蒙本的水循环、天蓝晚霞等课相互印证。

\subsection*{地理2 · 地貌、灾害与资源环境 \plan}

\textbf{对标：} 初中八上（地形+资源）+ 高中必修1（地貌）+ 高中选择性必修3精选\\
\textbf{内容：}
\begin{itemize}
  \item 海陆分布与板块构造：七大洲四大洋、板块构造学说
  \item 地形地貌：内力作用（褶皱/断层/火山）与外力作用（风化/侵蚀/搬运/堆积）、主要地貌类型（喀斯特/风蚀/冰川/海岸/河流）
  \item 土壤与植被：土壤的组成与剖面、植被与气候的对应关系
  \item 自然灾害：地震（成因+抗震）、火山喷发、滑坡与泥石流、洪涝与干旱、台风
  \item 自然资源：水资源、土地资源、矿产资源、能源的分类与分布
  \item 环境问题：全球变暖、臭氧层破坏、酸雨、荒漠化
  \item 地理技术入门：遥感（卫星影像判读）、GIS（图层叠加概念）、GPS/北斗
\end{itemize}
\textbf{定位：} 从"是什么"过渡到"为什么"和"怎么办"。自然灾害部分可呼应启蒙本的地震课和恐龙灭绝课。地理技术（GIS/遥感）是未来的基础工具，鼓励用Google Earth做探索。

\newpage

% ============================================================
\section*{\textcolor{orange!80!black}{六、课外拓展系列（3本·大纲占位）}}
% ============================================================

\textcolor{gray}{注：拓展系列仅列大纲骨架，不做详细内容展开。待基础系列完成后视需求启动。}

\subsection*{拓展1 · 物理竞赛思维与实验 \plan}

\textbf{面向：} 对物理有浓厚兴趣或竞赛需求\\
\textbf{框架：}
\begin{enumerate}
  \item 力学专题深化：复杂受力分析、连接体问题、质心与质心运动
  \item 电磁学专题深化：复杂电路等效变换、带电粒子在复合场中的运动
  \item 实验方法与误差分析：系统误差/偶然误差、控制变量法、图像法
  \item 开放性实验设计：给定器材设计方案、生活场景中的物理量测量
  \item 物理建模与估算：实际问题$\rightarrow$物理模型简化
  \item 竞赛入门：全国初中/高中物理竞赛题型与策略
\end{enumerate}

\subsection*{拓展2 · 化生实验与科学探究 \plan}

\textbf{面向：} 对化学或生物实验感兴趣\\
\textbf{框架：}
\begin{enumerate}
  \item 实验安全与基本操作规范
  \item 家庭化学实验：自制指示剂、晶体培养、自制肥皂、电解水
  \item 物质的检验与鉴别：常见气体/离子的检验
  \item 化学定量实验：溶液配制、酸碱中和滴定
  \item 显微镜实验：临时装片制作、微生物观察
  \item 生物探究项目：种子萌发条件、水果催熟、水质生物检测
  \item 科学论文写作：选题$\rightarrow$实验设计$\rightarrow$数据分析$\rightarrow$结论
  \item 竞赛入门：初中化学/生物竞赛题型特点
\end{enumerate}

\subsection*{拓展3 · 地球与环境科学探究 \plan}

\textbf{面向：} 对地理和环境科学有探究兴趣\\
\textbf{框架：}
\begin{enumerate}
  \item 气象观测实践：自制雨量器/风向标、天气数据分析
  \item 地图与GIS实操：Google Earth探索、简易地图绘制
  \item 野外考察指南：岩石矿物识别、土壤类型判断、河流地貌观察
  \item 北京在地项目：香山物候观测、永定河河漫滩考察、城市热岛效应测量
  \item 环境问题案例：密云水库水质、北京空气质量年际变化
  \item 可持续发展方案设计：校园节水/低碳出行规划
  \item 竞赛入门：全国地理奥赛（IGU）题型特点、野外考察题模拟
\end{enumerate}

\newpage

% ============================================================
\section*{\textcolor{blue!70!black}{附录：统计汇总}}
% ============================================================

\begin{center}
\begin{tabular}{|l|c|c|c|}
\hline
\textbf{系列} & \textbf{本数} & \textbf{状态} & \textbf{覆盖学段} \\
\hline
启蒙 & 1 & \done & 小学四--六年级 \\
\hline
物理 & 6 & \plan & 初中物理 + 高中物理（含选必） \\
化学 & 4 & \plan & 初中化学 + 高中化学（含选必） \\
生物 & 4 & \plan & 初中生物 + 高中生物（含选必） \\
地理 & 2 & \plan & 初中自然地理 + 高中自然地理（含选必精选） \\
\hline
拓展 & 3 & \plan（占位） & 跨学段（竞赛/探究导向） \\
\hline
\textbf{合计} & \textbf{20} & 1完成 / 19待写 & \\
\hline
\end{tabular}
\end{center}

\vspace{0.3cm}

\noindent\textbf{各系列推荐阅读顺序：}
\begin{itemize}
  \item 物理：1（力学入门）$\rightarrow$ 2（声光热）$\rightarrow$ 3（电磁入门）$\rightarrow$ 4（力学与运动定律）$\rightarrow$ 5（曲线、能量与动量）$\rightarrow$ 6（电磁、波与近代物理）
  \item 化学：1（化学语言）$\rightarrow$ 2（元素与反应）$\rightarrow$ 3（反应原理）$\rightarrow$ 4（有机与结构）
  \item 生物：1（分子与细胞）$\rightarrow$ 2（结构层次）$\rightarrow$ 3（遗传与进化）$\rightarrow$ 4（生态与多样性）
  \item 地理：1（自然地理基础）$\rightarrow$ 2（地貌、灾害与资源环境）
  \item 拓展：建议完成对应学科的基础本后再选读
\end{itemize}

\vspace{0.5cm}

\begin{center}
\rule{0.5\textwidth}{0.5pt}
\vspace{0.3cm}
\textcolor{gray}{\small 大纲版本：2026-05-22 · 审核通过}\\
\textcolor{gray}{\small 覆盖北京中考（物理计分；化/生/地等级考查）及高考3+3选考要求}\\
\textcolor{gray}{\small 人文/经济/区域地理内容归入 social-science/ 系列}
\end{center}

\end{document}
